\scanpage{250}
\mspage{200}

Quant au foncteur envisagé tantôt

\[
\unsure{\mathcal C|\mathscr C}\longrightarrow
\underbrace{\underline{\mathrm{Hom}}(\mathrm{Cat}_{/S},\underline{\mathrm{Hom}}(S^{\unsure{\circ|o}},\underline{\mathrm{Hot}}_W\deleted{\unsure{)}}))}_{\simeq\mathrm{Hom}(\mathrm{Cat}_{/S}\times S^\circ,\mathrm{Hot}_W)}
\]

\[
C\longmapsto ((X,s)\longmapsto\unsure{\mathrm{hot}|\mathrm{Hot}}_W(C(X)_s)),
\]

il se déduit du précédent en composant avec

\diagram{\begin{tikzcd}[ampersand replacement=\&,row sep=.6em,column sep=1em]
\mathrm{HOT}_W(S)\arrow[rr]\arrow[dr,dashed]\&\&\underline{\mathrm{Hom}}(S^\circ,\mathrm{Hot}_W)\\
\&\underline{\mathrm{Hom}}(\Pi S^\circ,\mathrm{Hot}_W)\arrow[ur,dashed]\&
\end{tikzcd}}

plus exactement avec le foncteur qui s'en déduit

\[
\underline{\mathrm{Hom}}(\mathrm{Cat}_{/S},\mathrm{HOT}_W(S))\longrightarrow\underline{\mathrm{Hom}}(\mathrm{Cat}_{/S},\underline{\mathrm{Hom}}(S^\circ,\mathrm{Hot}_W)\unsure{)|,}
\]

On voit poindre ici un principe plausible de canonicité \inserted{d'isomorphismes}, en ce qui concerne les foncteurs

\[
C_W:\mathrm{Cat}_{/S}\xrightarrow{\unsure{\leftarrow|\sim}}\mathrm{HOT}^{\unsure{\mathrm{lc}|\mathrm{le}}}_W(S)
\]

définis par les $C\in\mathrm{Ob}\,\mathcal C$, par la relation

\[
\begin{aligned}
\Pi\mathcal C&\approx\text{catégorie ponctuelle}\\
&\text{i.e.}\qquad\mathcal C\ \text{est }1\text{-connexe}.
\end{aligned}
\]

De plus, les arguments utilisés jusqu'ici sont de nature à tel point formelle, que les aspects spécifiques du Modéliseur $\mathrm{Cat}$
% archive sheet number: 250
